
\clearpage
\pdfbookmark{Abstract}{Abstract}
\section*{Abstract}

\begin{center}
\begin{minipage}[t]{0.7\textwidth}
	Implementations of programming languages are expected to faithfully realize their specifications, yet in practice this correspondence is often validated only by testing against written standards or reference implementations. 
	Since testing ranges over finitely many programs, it cannot establish the absence of mismatches for all well-typed inputs. 
	This thesis develops a machine-checked approach in Agda in which an implementation is systematically derived from its formal semantics and proved correct by construction. 
	We present a uniform derivation of abstract machines from semantic definitions, carried out for three case studies: arithmetic expressions, arithmetic expressions with exceptions, and the simply typed lambda calculus. 
	The derivation makes evaluation structure explicit via typed frames and a stack of machine states, with operational modes for evaluation and return, and an additional unwind mode for exceptions. 
	For each language, transition rules are obtained by mimicking in-order evaluation, and the resulting machine is proved correct and complete with respect to big-step semantics. 
	As a corollary, we establish semantic equivalences between big-step, small-step, denotational, and abstract-machine presentations, showing that the semantics are compositional and mutually consistent. 
	These results provide a systematic methodology for deriving verified interpreters from formal specifications.
\end{minipage}
\end{center}

\vfill

\section*{Zusammenfassung}

\begin{center}
\begin{minipage}[t]{0.7\textwidth}
	Von Implementierungen von Programmiersprachen wird erwartet, dass sie ihre Spezifikationen treu umsetzen; in der Praxis wird diese Übereinstimmung jedoch oft nur durch Tests anhand von Spezifikationsdokumenten 
	oder Referenzimplementierungen überprüft. Da sich die Tests auf eine endliche Anzahl von Programmen beschränken, können sie nicht für alle wohlgetypten Eingaben das Fehlen von Abweichungen nachweisen.
	Diese Arbeit entwickelt einen maschinell überprüfbaren Ansatz in Agda, bei dem eine Implementierung systematisch aus ihrer formalen Semantik abgeleitet und durch Konstruktion als korrekt bewiesen wird.
	Wir präsentieren eine einheitliche Ableitung abstrakter Maschinen aus semantischen Definitionen, die für drei Fallstudien durchgeführt wurde: arithmetische Ausdrücke, arithmetische Ausdrücke mit Ausnahmen und der einfach typisierte Lambda-Kalkül.
	Die Ableitung macht die Auswertungsstruktur durch typisierte Frames und einen Stapel von Maschinenzuständen explizit, mit Betriebsmodi für die Auswertung und die Rückgabe sowie einem zusätzlichen Unwind-Modus für Ausnahmen.
	Für jede Sprache werden Übergangsregeln durch Nachahmung der In-Order-Auswertung gewonnen, und die resultierende Maschine wird in Bezug auf die Big-Step-Semantik als korrekt und vollständig bewiesen.
	Als Korollar stellen wir semantische Äquivalenzen zwischen Big-Step-, Small-Step-, denotationalen und abstrakten Maschinen-Darstellungen her und zeigen, dass die Semantiken kompositorisch und wechselseitig konsistent sind.
	Diese Ergebnisse liefern eine systematische Methodik zur Ableitung verifizierter Interpreter aus formalen Spezifikationen.
\end{minipage}
\end{center}

\vfill

\thispagestyle{empty}